% Clase 12 — el paso clave de la prueba de phi_i = sum_j P_ij Q_j:
% "prender" solo una carga. El conductor 1 cargado, los demas neutros; los
% potenciales que aparecen se rotulan con el superindice de la carga prendida.
\begin{tikzpicture}[scale=1]

  % conductor 1 (cargado)
  \fill[figred!10, draw=figblue, line width=0.9pt]
    plot[smooth cycle, tension=0.8] coordinates {(-0.9,0.2) (-0.2,1.0) (0.8,0.6) (0.6,-0.5) (-0.5,-0.6)};
  \foreach \p in {(-0.75,0.25),(-0.2,0.85),(0.6,0.55),(0.5,-0.35),(-0.4,-0.45)}{
    \node[figlbl, figred] at \p {$+$};}
  \node[figlbl] at (0,0.15) {$1$};
  \node[figlbl] at (0,-1.15) {$Q_1=\bar Q_1$};
  \node[figlblaux] at (0,-1.6) {$\phi_1^{(1)}$};

  % conductor 2 (neutro)
  \fill[figgray!25, draw=figblue, line width=0.9pt]
    plot[smooth cycle, tension=0.8] coordinates {(3.1,1.4) (3.9,2.0) (4.6,1.4) (4.0,0.8)};
  \node[figlbl] at (3.9,1.4) {$2$};
  \node[figlbl] at (5.5,1.75) {$Q_2=0$};
  \node[figlblaux] at (5.5,1.3) {$\phi_2^{(1)}$};

  % conductor 3 (neutro)
  \fill[figgray!25, draw=figblue, line width=0.9pt]
    plot[smooth cycle, tension=0.8] coordinates {(3.0,-1.2) (3.8,-0.7) (4.7,-1.1) (4.3,-1.9) (3.3,-1.9)};
  \node[figlbl] at (3.85,-1.3) {$3$};
  \node[figlbl] at (5.7,-1.05) {$Q_3=0$};
  \node[figlblaux] at (5.7,-1.5) {$\phi_3^{(1)}$};

\end{tikzpicture}
